\documentclass{stex}

\libusepackage{naproche}
\libinput{preamble}

\begin{document}
\begin{smodule}{computation-laws-for-classes.ftl}

\importmodule[libraries/foundations]{classes.ftl}

\begin{proposition*}[forthel]
  Let $A, B$ be classes.
  Then $A \Union B \Eq B \Union A$.
\end{proposition*}
\begin{proof}[forthel]
  Let us show that $A \Union B \Subclass B \Union A$.
    Let $x \Elem A \Union B$.
    Then $x \Elem A$ or $x \Elem B$.
    Hence $x \Elem B$ or $x \Elem A$.
    Thus $x \Elem B \Union A$.
  End.

  Let us show that $B \Union A \Subclass A \Union B$.
    Let $x \Elem B \Union A$.
    Then $x \Elem B$ or $x \Elem A$.
    Hence $x \Elem A$ or $x \Elem B$.
    Thus $x \Elem A \Union B$.
  End.
\end{proof}

\begin{proposition*}[forthel]
  Let $A, B$ be classes.
  Then $A \Intersect B \Eq B \Intersect A$.
\end{proposition*}
\begin{proof}[forthel]
  Let us show that $A \Intersect B \Subclass B \Intersect A$.
    Let $x \Elem A \Intersect B$.
    Then $x \Elem A$ and $x \Elem B$.
    Hence $x \Elem B$ and $x \Elem A$.
    Thus $x \Elem B \Intersect A$.
  End.

  Let us show that $B \Intersect A \Subclass A \Intersect B$.
    Let $x \Elem B \Intersect A$.
    Then $x \Elem B$ and $x \Elem A$.
    Hence $x \Elem A$ and $x \Elem B$.
    Thus $x \Elem A \Intersect B$.
  End.
\end{proof}

\begin{proposition*}[forthel]
  Let $A, B, C$ be classes.
  Then $(A \Union B) \Union C \Eq A \Union (B \Union C)$.
\end{proposition*}
\begin{proof}[forthel]
  Let us show that $((A \Union B) \Union C) \Subclass A \Union (B \Union C)$. %!
    Let $x \Elem (A \Union B) \Union C$.
    Then $x \Elem A \Union B$ or $x \Elem C$.
    Hence $x \Elem A$ or $x \Elem B$ or $x \Elem C$.
    Thus $x \Elem A$ or $x \Elem (B \Union C)$.
    Therefore $x \Elem A \Union (B \Union C)$.
  End.

  Let us show that $A \Union (B \Union C) \Subclass (A \Union B) \Union C$.
    Let $x \Elem A \Union (B \Union C)$.
    Then $x \Elem A$ or $x \Elem B \Union C$.
    Hence $x \Elem A$ or $x \Elem B$ or $x \Elem C$.
    Thus $x \Elem A \Union B$ or $x \Elem C$.
    Therefore $x \Elem (A \Union B) \Union C$.
  End.
\end{proof}

\begin{proposition*}[forthel]
  Let $A, B, C$ be classes.
  Then $(A \Intersect B) \Intersect C \Eq A \Intersect (B \Intersect C)$.
\end{proposition*}
\begin{proof}[forthel]
  Let us show that $((A \Intersect B) \Intersect C) \Subclass A \Intersect (B \Intersect C)$. %!
    Let $x \Elem (A \Intersect B) \Intersect C$.
    Then $x \Elem A \Intersect B$ and $x \Elem C$.
    Hence $x \Elem A$ and $x \Elem B$ and $x \Elem C$.
    Thus $x \Elem A$ and $x \Elem (B \Intersect C)$.
    Therefore$ x \Elem A \Intersect (B \Intersect C)$.
  End.

  Let us show that $A \Intersect (B \Intersect C) \Subclass (A \Intersect B) \Intersect C$.
    Let $x \Elem A \Intersect (B \Intersect C)$.
    Then $x \Elem A$ and $x \Elem B \Intersect C$.
    Hence $x \Elem A$ and $x \Elem B$ and $x \Elem C$.
    Thus $x \Elem A \Intersect B$ and $x \Elem C$.
    Therefore $x \Elem (A \Intersect B) \Intersect C$.
  End.
\end{proof}

\begin{proposition*}[forthel]
  Let $A, B, C$ be classes.
  Then $A \Intersect (B \Union C) \Eq (A \Intersect B) \Union (A \Intersect C)$.
\end{proposition*}
\begin{proof}[forthel]
  Let us show that $A \Intersect (B \Union C) \Subclass (A \Intersect B) \Union (A \Intersect C)$.
    Let $x \Elem A \Intersect (B \Union C)$.
    Then $x \Elem A$ and $x \Elem B \Union C$.
    Hence $x \Elem A$ and ($x \Elem B$ or $x \Elem C$).
    Thus ($x \Elem A$ and $x \Elem B$) or ($x \Elem A$ and $x \Elem C$).
    Therefore $x \Elem A \Intersect B$ or $x \Elem A \Intersect C$.
    Hence $x \Elem (A \Intersect B) \Union (A \Intersect C)$.
  End.

  Let us show that $((A \Intersect B) \Union (A \Intersect C)) \Subclass A \Intersect (B \Union C)$. %!
    Let $x \Elem (A \Intersect B) \Union (A \Intersect C)$.
    Then $x \Elem A \Intersect B$ or $x \Elem A \Intersect C$.
    Hence ($x \Elem A$ and $x \Elem B$) or ($x \Elem A$ and $x \Elem C$).
    Thus $x \Elem A$ and ($x \Elem B$ or $x \Elem C$).
    Therefore $x \Elem A$ and $x \Elem B \Union C$.
    Hence$ x \Elem A \Intersect (B \Union C)$.
  End.
\end{proof}

\begin{proposition*}[forthel]
  Let $A, B, C$ be classes.
  Then $A \Union (B \Intersect C) \Eq (A \Union B) \Intersect (A \Union C)$.
\end{proposition*}
\begin{proof}[forthel]
  Let us show that $A \Union (B \Intersect C) \Subclass (A \Union B) \Intersect (A \Union C)$.
    Let $x \Elem A \Union (B \Intersect C)$.
    Then $x \Elem A$ or $x \Elem B \Intersect C$.
    Hence $x \Elem A$ or ($x \Elem B$ and $x \Elem C$).
    Thus ($x \Elem A$ or $x \Elem B$) and ($x \Elem$ A or $x \Elem C$).
    Therefore $x \Elem A \Union B$ and $x \Elem A \Union C$.
    Hence $x \Elem (A \Union B) \Intersect (A \Union C)$.
  End.

  Let us show that $((A \Union B) \Intersect (A \Union C)) \Subclass A \Union (B \Intersect C)$. %!
    Let $x \Elem (A \Union B) \Intersect (A \Union C)$.
    Then $x \Elem A \Union B$ and $x \Elem A \Union C$.
    Hence ($x \Elem A$ or $x \Elem B$) and ($x \Elem A$ or $x \Elem C$).
    Thus $x \Elem A$ or ($x \Elem B$ and $x \Elem C$).
    Therefore $x \Elem A$ or $x \Elem B \Intersect C$.
    Hence $x \Elem A \Union (B \Intersect C)$.
  End.
\end{proof}

\begin{proposition*}[forthel]
  Let $A$ be a class.
  Then $A \Union A \Eq A$.
\end{proposition*}
\begin{proof}[forthel]
  $A \Union A \Eq \CClass{x}{x \Elem A\text{ or }x \Elem A}$.
  Hence $A \Union A \Eq \CClass{x}{x \Elem A}$.
  Thus $A \Union A \Eq A$.
\end{proof}

\begin{proposition*}[forthel]
  Let $A$ be a class.
  Then $A \Intersect A \Eq A$.
\end{proposition*}
\begin{proof}[forthel]
  $A \Intersect A \Eq \CClass{x}{x \Elem A\text{ and }x \Elem A}$.
  Hence $A \Intersect A \Eq \CClass{x}{x \Elem A}$.
  Thus $A \Intersect A \Eq A$.
\end{proof}

\begin{proposition*}[forthel]
  Let $A, B, C$ be classes.
  Then $A \Diff (B \Intersect C) \Eq (A \Diff B) \Union (A \Diff C)$.
\end{proposition*}
\begin{proof}[forthel]
  Let us show that $A \Diff (B \Intersect C) \Subclass (A \Diff B) \Union (A \Diff C)$.
    Let $x \Elem A \Diff (B \Intersect C)$.
    Then $x \Elem A$ and $x \NotElem B \Intersect C$.
    Hence it is wrong that ($x \Elem B$ and $x \Elem C$).
    Thus $x \NotElem B$ or $x \NotElem C$.
    Therefore $x \Elem A$ and ($x \NotElem B$ or $x \NotElem C$).
    Then ($x \Elem A$ and $x \NotElem B$) or ($x \Elem A$ and $x \NotElem C$).
    Hence $x \Elem A \Diff B$ or $x \Elem A \Diff C$.
    Thus $x \Elem (A \Diff B) \Union (A \Diff C)$.
  End.

  Let us show that $((A \Diff B) \Union (A \Diff C)) \Subclass A \Diff (B \Intersect C)$. %!
    Let $x \Elem (A \Diff B) \Union (A \Diff C)$.
    Then $x \Elem A \Diff B$ or $x \Elem A \Diff C$.
    Hence ($x \Elem A$ and $x \NotElem B$) or ($x \Elem A$ and $x \NotElem C$).
    Thus $x \Elem A$ and ($x \NotElem B$ or $x \NotElem C$).
    Therefore $x \Elem A$ and $\Not(x \Elem B \And x \Elem C)$.
    Then $x \Elem A$ and $x \NotElem B \Intersect C$.
    Hence $x \Elem A \Diff (B \Intersect C)$.
  End.
\end{proof}

\begin{proposition*}[forthel]
  Let $A, B, C$ be classes.
  Then $A \Diff (B \Union C) \Eq (A \Diff B) \Intersect (A \Diff C)$.
\end{proposition*}
\begin{proof}[forthel]
  Let us show that $A \Diff (B \Union C) \Subclass (A \Diff B) \Intersect (A \Diff C)$.
    Let $x \Elem A \Diff (B \Union C)$.
    Then $x \Elem A$ and $x \NotElem B \Union C$.
    Hence it is wrong that ($x \Elem B$ or $x \Elem C$).
    Thus $x \NotElem B$ and $x \NotElem C$.
    Therefore $x \Elem A$ and ($x \NotElem B$ and $x \NotElem C$).
    Then ($x \Elem A$ and $x \NotElem B$) and ($x \Elem A$ and $x \NotElem C$).
    Hence $x \Elem A \Diff B$ and $x \Elem A \Diff C$.
    Thus $x \Elem (A \Diff B) \Intersect (A \Diff C)$.
  End.

  Let us show that $((A \Diff B) \Intersect (A \Diff C)) \Subclass A \Diff (B \Union C)$. %!
    Let $x \Elem (A \Diff B) \Intersect (A \Diff C)$.
    Then $x \Elem A \Diff B$ and $x \Elem A \Diff C$.
    Hence ($x \Elem A$ and $x \NotElem B$) and ($x \Elem A$ and $x \NotElem C$).
    Thus $x \Elem A$ and ($x \NotElem B$ and $x \NotElem C$).
    Therefore $x \Elem A$ and $\Not(x \Elem B \Or x \Elem C)$.
    Then $x \Elem A$ and $x \NotElem B \Union C$.
    Hence $x \Elem A \Diff (B \Union C)$.
  End.
\end{proof}

\begin{proposition*}[forthel]
  Let $A, B$ be classes.
  Then $A \Subclass A \Union B$.
\end{proposition*}
\begin{proof}[forthel]
  Let $x \Elem A$.
  Then $x \Elem A$ or $x \Elem B$.
  Hence $x \Elem A \Union B$.
\end{proof}

\begin{proposition*}[forthel]
  Let $A, B$ be classes.
  Then $A \Intersect B \Subclass A$.
\end{proposition*}
\begin{proof}[forthel]
  Let $x \Elem A \Intersect B$.
  Then $x \Elem A$ and $x \Elem B$.
  Hence $x \Elem A$.
\end{proof}

\begin{proposition*}[forthel]
  Let $A, B$ be classes.
  Then $A \Subclass B$ iff $A \Union B \Eq B$.
\end{proposition*}
\begin{proof}[forthel]
  \begin{case}{$A \Subclass B$.}

    Let us show that $A \Union B \Subclass B$.
      Let $x \Elem A \Union B$.
      Then $x \Elem A$ or $x \Elem B$.
      If $x \Elem A$ then $x \Elem B$.
      Hence $x \Elem B$.
    End.

    Let us show that $B \Subclass A \Union B$.
      Let $x \Elem B$.
      Then $x \Elem A$ or $x \Elem B$.
      Hence $x \Elem A \Union B$.
    End.
  \end{case}

  \begin{case}{$A \Union B \Eq B$.}
    Let $x \Elem A$.
    Then $x \Elem A$ or $x \Elem B$.
    Hence $x \Elem A \Union B \Eq B$.
  \end{case}
\end{proof}

\begin{proposition*}[forthel]
  Let $A, B$ be classes.
  Then $A \Subclass B$ iff $A \Intersect B \Eq A$.
\end{proposition*}
\begin{proof}[forthel]
  \begin{case}{$A \Subclass B$.}

    Let us show that $A \Intersect B \Subclass A$.
      Let $x \Elem A \Intersect B$.
      Then $x \Elem A$ and $x \Elem B$.
      Hence $x \Elem A$.
    End.

    Let us show that $A \Subclass A \Intersect B$.
      Let $x \Elem A$.
      Then $x \Elem B$.
      Hence $x \Elem A$ and $x \Elem B$.
      Thus $x \Elem A \Intersect B$.
    End.
  \end{case}

  \begin{case}{$A \Intersect B \Eq A$.}
    Let $x \Elem A$.
    Then $x \Elem A \Intersect B$.
    Hence $x \Elem A$ and $x \Elem B$.
    Thus $x \Elem B$.
  \end{case}
\end{proof}

\begin{proposition*}[forthel]
  Let $A$ be a class.
  Then $A \Diff A \Eq \EmptyClass$.
\end{proposition*}
\begin{proof}[forthel]
  $A \Diff A$ has no elements.
  Indeed $A \Diff A \Eq \CClass{x}{x \Elem A\text{ and }x \NotElem A}$.
  Hence the thesis.
\end{proof}

\begin{proposition*}[forthel]
  Let $A$ be a class.
  Then $A \Diff \EmptyClass \Eq A$.
\end{proposition*}
\begin{proof}[forthel]
  $A \Diff \EmptyClass \Eq \CClass{x}{x \Elem A\text{ and }x \NotElem \EmptyClass}$.
  No element is an element of $\EmptyClass$.
  Hence $A \Diff \EmptyClass \Eq \CClass{x}{x \Elem A}$.
  Then we have the thesis.
\end{proof}

\begin{proposition*}[forthel]
  Let $A, B$ be classes.
  Then $A \Diff (A \Diff B) \Eq A \Intersect B$.
\end{proposition*}
\begin{proof}[forthel]
  Let us show that $A \Diff (A \Diff B) \Subclass A \Intersect B$.
    Let $x \Elem A \Diff (A \Diff B)$.
    Then $x \Elem A$ and $x \NotElem A \Diff B$.
    Hence $x \NotElem A$ or $x \Elem B$.
    Thus $x \Elem B$.
    Therefore $x \Elem A \Intersect B$.
  End.

  Let us show that $A \Intersect B \Subclass A \Diff (A \Diff B)$.
    Let $x \Elem A \Intersect B$.
    Then $x \Elem A$ and $x \Elem B$.
    Hence $x \NotElem A$ or $x \Elem B$.
    Thus $x \NotElem A \Diff B$.
    Therefore $x \Elem A \Diff (A \Diff B)$.
  End.
\end{proof}

\begin{proposition*}[forthel]
  Let $A, B$ be classes.
  Then $B \Subclass A$ iff $A \Diff (A \Diff B) \Eq B$.
\end{proposition*}
\begin{proof}[forthel]
  \begin{case}{$B \Subclass A$.} \end{case}

  \begin{case}{$A \Diff (A \Diff B) \Eq B$.}
    Then every element of $B$ is an element of $A \Diff (A \Diff B)$.
    Thus every element of $B$ is an element of $A$.
    Then we have the thesis.
  \end{case}
\end{proof}

\begin{proposition*}[forthel]
  Let $A, B, C$ be classes.
  Then $A \Intersect (B \Diff C) \Eq (A \Intersect B) \Diff (A \Intersect C)$.
\end{proposition*}
\begin{proof}[forthel]
  Let us show that $A \Intersect (B \Diff C) \Subclass (A \Intersect B) \Diff (A \Intersect C)$.
    Let $x \Elem A \Intersect (B \Diff C)$.
    Then $x \Elem A$ and $x \Elem B \Diff C$.
    Hence $x \Elem A$ and $x \Elem B$.
    Thus $x \Elem A \Intersect B$ and $x \NotElem C$.
    Therefore $x \NotElem A \Intersect C$.
    Then we have $x \Elem (A \Intersect B) \Diff (A \Intersect C)$.
  End.

  Let us show that $((A \Intersect B) \Diff (A \Intersect C)) \Subclass A \Intersect (B \Diff C)$. %!
    Let $x \Elem (A \Intersect B) \Diff (A \Intersect C)$.
    Then $x \Elem A$ and $x \Elem B$.
    $x \NotElem A \Intersect C$.
    Hence $x \NotElem C$.
    Thus $x \Elem B \Diff C$.
    Therefore $x \Elem A \Intersect (B \Diff C)$.
  End.
\end{proof}

\end{smodule}
\end{document}
